\section*{Bab \ref{ch:b1:intermolecular-forces-solvents} --- Gaya Antarmolekul dan Pelarut}

\begin{solution}{exo:b1:intermolecular-forces-solvents:1}
Argon: hanya London. Hidrogen klorida: Keesom, Debye, dan London, dengan
London yang dominan (\cref{exo:b1:intermolecular-forces-solvents:8});
klor terlalu besar untuk membentuk \omterm{def:b1:intermolecular-forces-solvents:hbond}{ikatan hidrogen} yang kuat. Air: \omterm{def:b1:intermolecular-forces-solvents:hbond}{ikatan
hidrogen} dominan, disertai ketiga suku van der Waals. Metana: hanya
London. Diiodin: hanya London, yang kuat karena molekulnya besar dan
sangat mudah terpolarisasi.
\end{solution}

\begin{solution}{exo:b1:intermolecular-forces-solvents:2}
Di antara molekul-molekulnya sendiri: \ce{CH3OH} (\ce{O-H}), \ce{CH3NH2}
(\ce{N-H}), dan \ce{HF}. \ce{CH3OCH3} dan \ce{CH3F} mempunyai \omterm{def:b1:lewis-resonance-vsepr:lewis}{pasangan
elektron bebas} pada O atau F, tetapi tidak ada hidrogen yang terikat
padanya: keduanya hanya dapat menerima \omterm{def:b1:intermolecular-forces-solvents:hbond}{ikatan hidrogen}, misalnya dari
air. \ce{CH4} tidak memberi dan tidak menerima.
\end{solution}

\begin{solution}{exo:b1:intermolecular-forces-solvents:3}
Air dan etanol: polar, protik. Propanon dan asetonitril: polar, aprotik.
Diklorometana: agak polar, aprotik. Heksana: apolar, aprotik.
\end{solution}

\begin{solution}{exo:b1:intermolecular-forces-solvents:4}
Diiodin apolar: di dalam heksana ia menukar \omterm{def:b1:intermolecular-forces-solvents:vdw}{interaksi London} dengan
\omterm{def:b1:intermolecular-forces-solvents:vdw}{interaksi London}, sedangkan di dalam air ia harus memutus \omterm{def:b1:intermolecular-forces-solvents:hbond}{ikatan hidrogen}
tanpa menggantinya. Natrium klorida bersifat ionik: hanya pelarut
berpermitivitas tinggi yang dapat memisahkan ion-ionnya dan hanya
\omterm{def:b1:intermolecular-forces-solvents:solvent-classes}{pelarut polar} yang dapat menyolvasinya; heksana ($\varepsilon_r = 1.89$)
tidak dapat melakukan keduanya.
\end{solution}

\begin{solution}{exo:b1:intermolecular-forces-solvents:5}
Metoksimetana tidak mempunyai \ce{O-H}: zat ini tidak dapat memberi
\omterm{def:b1:intermolecular-forces-solvents:hbond}{ikatan hidrogen} dan mendidih paling rendah. Metanol dan etanol sama-sama
membentuk \omterm{def:b1:intermolecular-forces-solvents:hbond}{ikatan hidrogen}; etanol, dengan satu \ce{CH2} lebih banyak,
mempunyai \omterm{def:b1:intermolecular-forces-solvents:vdw}{interaksi London} yang lebih kuat. Asam etanoat membentuk dimer
berikatan hidrogen dan paling berat: zat ini mendidih paling tinggi.
\end{solution}

\begin{solution}{exo:b1:intermolecular-forces-solvents:6}
Di vakum $F = e^2/(4\pi\varepsilon_0 r^2) = (1.602 \times
10^{-19})^2/(4\pi \times 8.854 \times 10^{-12} \times (5.00 \times
10^{-10})^2) = \qty{9.23e-10}{N}$. Di dalam air $F/78.4 = \qty{1.18e-11}{N}$;
di dalam heksana $F/1.89 = \qty{4.88e-10}{N}$. Air melemahkan tarikan itu
41 kali lebih banyak daripada heksana: air dapat menjaga ion-ion tetap
terpisah, heksana tidak.
\end{solution}

\begin{solution}{exo:b1:intermolecular-forces-solvents:7}
Mencampur propanon dengan air memutus sebagian \omterm{def:b1:intermolecular-forces-solvents:hbond}{ikatan hidrogen} air--air,
tetapi membentuk \omterm{def:b1:intermolecular-forces-solvents:hbond}{ikatan hidrogen} baru antara air dan oksigen \ce{C=O},
ditambah interaksi dipol--dipol: neracanya menguntungkan. Heksana hanya
menawarkan \omterm{def:b1:intermolecular-forces-solvents:vdw}{interaksi London} yang lemah kepada air sebagai ganti \omterm{def:b1:intermolecular-forces-solvents:hbond}{ikatan
hidrogen} yang harus diputus: kedua zat cair itu tetap terpisah.
\end{solution}

\begin{solution}{exo:b1:intermolecular-forces-solvents:8}
\ce{HBr} mempunyai lebih banyak elektron dan awan elektron yang lebih
besar dan lebih mudah terpolarisasi: \omterm{def:b1:intermolecular-forces-solvents:vdw}{interaksi Londonnya} melampaui
\omterm{def:b1:intermolecular-forces-solvents:vdw}{interaksi London} \ce{HCl} lebih banyak daripada kekurangan \omterm{def:b1:intermolecular-forces-solvents:vdw}{interaksi
Keesomnya}. London yang dominan.
\end{solution}

\begin{solution}{exo:b1:intermolecular-forces-solvents:9}
\chemfig{CH_3-C(=[:60]O)-[:-60]O-[:0]H} berhadapan dengan pasangannya,
dua ikatan $\ce{O-H}\cdots\ce{O=C}$ menutup sebuah cincin beranggota
delapan. Dalam benzena, yang apolar dan aprotik, tidak ada yang menyaingi
ikatan-ikatan ini: asam itu terdimerisasi dan berperilaku sebagai partikel
bermassa sekitar \qty{120}{g/mol}. Dalam air, setiap molekul asam justru
membentuk \omterm{def:b1:intermolecular-forces-solvents:hbond}{ikatan hidrogen} dengan air: dimer-dimernya terurai.
\end{solution}

\begin{solution}{exo:b1:intermolecular-forces-solvents:10}
Natrium klorida tersusun atas ion-ion: air, yang permitivitasnya tinggi,
memisahkannya (disosiasi) dan menghidrasinya (oksigen ke arah \ce{Na+},
hidrogen ke arah \ce{Cl-}). Hidrogen klorida adalah \omterm{def:b1:lewis-resonance-vsepr:dipole}{molekul polar}: air,
yang polar, mengubah ikatan \ce{H-Cl} menjadi pasangan ion \ce{H3O+ Cl-}
(ionisasi), memisahkan pasangan itu (disosiasi), dan menghidrasi
ion-ionnya. Heksana apolar dan berpermitivitas rendah: pelarut ini tidak
mengionkan dan tidak mendisosiasikan, sehingga tidak ada ion yang
membawa arus.
\end{solution}

\begin{solution}{exo:b1:intermolecular-forces-solvents:11}
$(174.1 - 36.1)/5 = \qty{27.6}{\celsius}$ per \ce{CH2}. Pertambahan
terakhir (nonana ke dekana) adalah \qty{23.4}{\celsius}, sehingga undekana
semestinya mendidih di dekat $174 + 22 \approx \qty{196}{\celsius}$.
Setiap \ce{CH2} menambah sumbangan London yang sama, tetapi merupakan
bagian yang makin kecil dari keseluruhan molekul, sementara suhu yang
diperlukan untuk mengatasi interaksi itu sudah tinggi.
% ledger: bp:C9H20
\end{solution}

\begin{solution}{exo:b1:intermolecular-forces-solvents:12}
Bagian \omterm{def:b1:intermolecular-forces-solvents:amphiphile}{hidrofilik}: kepala sulfat \ce{-OSO3^-} dengan ion lawannya,
\ce{Na+}; bagian \omterm{def:b1:intermolecular-forces-solvents:amphiphile}{hidrofobik}: rantai dua belas karbon. Di permukaan,
kepala di dalam air dan ekor di udara; di dalam misel, ekor di dalam dan
kepala di luar. Ekor-ekornya larut dalam lemak noda, kepala-kepalanya
menjaga noda itu tetap bersentuhan dengan air, dan lemaknya terangkat
dalam bentuk misel.
\end{solution}

\begin{solution}{pb:b1:intermolecular-forces-solvents:1}
\textbf{1.} \omterm{def:b1:intermolecular-forces-solvents:vdw}{Interaksi London}, satu-satunya interaksi antara atom-atom
apolar.
\textbf{2.} Dari helium sampai xenon, banyaknya elektron dan ukurannya
bertambah, sehingga \omterm{def:b1:periodicity:polarisability}{polarisabilitas} dan tarikan London juga bertambah.
\textbf{3.} Helium mempunyai dua elektron yang diikat erat oleh intinya:
atom yang paling sukar terpolarisasi, dengan tarikan London yang paling
lemah.
\textbf{4.} $68.8 - 36.1 = \qty{32.7}{\celsius}$: satu \ce{CH2} menambah
elektron dan luas permukaan kontak, sehingga menambah tarikan London.
\textbf{5.} London, yang membesar dengan \omterm{def:b1:periodicity:polarisability}{polarisabilitas}: $\ce{CH4} < \ce{SiH4} <
\ce{GeH4}$.
\textbf{6.} Karbon kurang elektronegatif untuk memolarisasi \ce{C-H}
dengan kuat, dan metana tidak mempunyai \omterm{def:b1:lewis-resonance-vsepr:lewis}{pasangan elektron bebas} untuk
menerima \omterm{def:b1:intermolecular-forces-solvents:hbond}{ikatan hidrogen}.
\textbf{7.} Air: polar, protik, dengan permitivitas tertinggi; air
memisahkan dan menyolvasi kedua ion.
\textbf{8.} $78.4/1.89 = 41$: tarikannya 41 kali lebih lemah di dalam
air.
\textbf{9.} Asetonitril (atau propanon): cukup polar untuk melarutkan
garamnya, dan aprotik sehingga anionnya tetap bebas.
\textbf{10.} Heksana: tidak \omterm{def:b1:intermolecular-forces-solvents:miscibility}{dapat bercampur} dengan air, apolar seperti
diiodin. Iodin berpindah ke dalam heksana, lapisan atas (heksana kurang
rapat daripada air).
\textbf{11.} Etanol \omterm{def:b1:intermolecular-forces-solvents:miscibility}{dapat bercampur} dengan air: tidak terbentuk lapisan
kedua.
\textbf{12.} Oksigennya menerima \omterm{def:b1:intermolecular-forces-solvents:hbond}{ikatan hidrogen} dari air, tetapi
etoksietana tidak dapat memberi satu pun dan kedua gugus etilnya
\omterm{def:b1:intermolecular-forces-solvents:amphiphile}{hidrofobik}: zat ini hanya sedikit larut.
\textbf{13.} $-66.4 - (-85.0) = \qty{18.6}{\celsius}$ per \omterm{def:b1:periodicity:table}{periode};
\ce{HF} hasil ekstrapolasi: $-85.0 - 18.6 = \qty{-103.6}{\celsius}$.
\textbf{14.} \ce{HF} mendidih pada \qty{19.6}{\celsius},
\qty{123}{\celsius} di atas ekstrapolasi.
\textbf{15.} $-62.5 - (-87.8) = \qty{25.3}{\celsius}$; \ce{NH3} hasil
ekstrapolasi: \qty{-113.1}{\celsius}; nilai sebenarnya $-33.4$,
\qty{80}{\celsius} lebih tinggi.
\textbf{16.} $-88.1 - (-112.0) = \qty{23.9}{\celsius}$; \ce{CH4} hasil
ekstrapolasi: \qty{-135.9}{\celsius}; nilai sebenarnya $-161.5$, di
bawahnya. Metana tidak membentuk \omterm{def:b1:intermolecular-forces-solvents:hbond}{ikatan hidrogen}; molekulnya hanya kecil
dan sukar terpolarisasi.
\textbf{17.} Air (sekitar \qty{179}{\celsius}) $>$ \ce{HF} (123) $>$
\ce{NH3} (80).
\textbf{18.} \ce{NH3} dapat memberi tiga tetapi hanya menerima satu (satu
\omterm{def:b1:lewis-resonance-vsepr:lewis}{pasangan bebas}); \ce{HF} dapat menerima tiga tetapi hanya memberi satu.
Air memberi dua dan menerima dua: setiap molekul dapat terlibat dalam
empat \omterm{def:b1:intermolecular-forces-solvents:hbond}{ikatan hidrogen}, sebuah jejaring tiga dimensi yang lengkap.
\textbf{19.} \ce{H2S} dan \ce{H2Se} tidak membentuk \omterm{def:b1:intermolecular-forces-solvents:hbond}{ikatan hidrogen} yang
berarti: titik didihnya mengikuti kecenderungan London yang akan diikuti
air tanpa \omterm{def:b1:intermolecular-forces-solvents:hbond}{ikatan hidrogen}. Air sendiri adalah anomali yang sedang diukur.
\textbf{20.} $-41.3 - (-60.3) = \qty{19.0}{\celsius}$ per \omterm{def:b1:periodicity:table}{periode}.
\textbf{21.} $T(p) = -60.3 + 19.0\,(p - 3)$ dalam \unit{\celsius}.
\textbf{22.} $T(2) = -60.3 - 19.0 = \qty{-79.3}{\celsius}$.
\textbf{23.} $100.0 - (-79.3) = \qty{179}{\celsius}$.
\textbf{24.} Tanpa \omterm{def:b1:intermolecular-forces-solvents:hbond}{ikatan hidrogennya} air akan mendidih di dekat
\textbf{\qty{-79}{\celsius}}, dan akan berwujud gas di Bumi.
\end{solution}
